\documentclass[11pt]{article}

\newcommand{\MemoTitle}{The Next AI Bottleneck Is Delegated Authority}
\newcommand{\MemoDate}{April 18, 2026}
\newcommand{\MemoLogoPath}{../../../tex/assets/icon-180.png}

% =========================================================
% Shared LaTeX preamble for all memos/*/tex/memo.tex
% =========================================================

\usepackage[letterpaper,top=1in,bottom=1in,left=1in,right=1in]{geometry}
\usepackage{microtype}
\usepackage[T1]{fontenc}
\usepackage[utf8]{inputenc}
\usepackage{lmodern}
\usepackage{graphicx}
\usepackage{xcolor}
\usepackage{hyperref}
\usepackage{xurl}
\usepackage{enumitem}

\setlength{\parindent}{0pt}
\setlength{\parskip}{0.35em}
\setlist[itemize]{leftmargin=1.2em,itemsep=0.04em,topsep=0.08em}
\setlist[enumerate]{leftmargin=1.4em,itemsep=0.04em,topsep=0.08em}

\definecolor{accentgold}{HTML}{D4A574}
\definecolor{ink}{HTML}{0F172A}

\hypersetup{
  hidelinks,
  pdftitle={\MemoTitle},
  pdfauthor={Adam Boas},
  pdfsubject={Policy Memo}
}

\pagestyle{empty}
\sloppy

\providecommand{\MemoTitle}{Memo}
\providecommand{\MemoDate}{\today}
\providecommand{\MemoLogoPath}{../../../tex/assets/icon-180.png}
\providecommand{\MemoAuthorName}{Adam Boas}
\providecommand{\MemoAuthorSite}{https://adamboas.com/}
\providecommand{\MemoAuthorSiteLabel}{adamboas.com}
\providecommand{\MemoAuthorEmail}{anboas@gmail.com}

\providecommand{\tightlist}{}

\newcommand{\MemoRenderHeader}{%
  \vspace*{-3.15em}
  \noindent
  \begin{minipage}[c]{0.14\textwidth}
  \raggedright
  \includegraphics[width=0.66\linewidth]{\MemoLogoPath}
  \end{minipage}%
  \begin{minipage}[c]{0.86\textwidth}
  \centering
  {\Large\bfseries\textcolor{ink}{\MemoAuthorName}}\\[-0.1em]
  {\small\textcolor{ink}{\href{\MemoAuthorSite}{\MemoAuthorSiteLabel} \quad|\quad \href{mailto:\MemoAuthorEmail}{\MemoAuthorEmail}}}\\[-0.1em]
  {\small\textcolor{ink}{\MemoDate}}
  \end{minipage}

  \vspace{0.35em}
  {\centering\large\bfseries \MemoTitle\par}
  \vspace{0.28em}
  \textcolor{accentgold}{\rule{\textwidth}{1.2pt}}
}


\begin{document}
\MemoRenderHeader

Everyone is still arguing about models.

Government is starting to ask a harder question: what does it take to
trust software that can act?

That is the real signal in the current stack. NIST is pushing on agent
standards, agent security, identity and authorization, evaluation,
monitoring, and procurement-facing measurement. DARPA is pushing on
trustworthy systems, AI engineering, and human-AI teaming. Defense
acquisition is starting to ask for open implementations, public
documentation, conformance artifacts, and software that other teams can
inherit.

These are not the same effort.

They are, however, circling the same missing layer.

The next bottleneck in AI is not generation.

It is governed delegated authority.

That has been the through-line of the argument on this site for a while
now.

\emph{From AI Force Multiplication to Force Creation} argued that the
strategic shift is not from no AI to some AI. It is from human-hours
amplified to agent-hours delegated. \emph{ACP-RA} argued that this
transition is not fieldable without a control plane: explicit identities
for agents, explicit trust scopes, mediated tool use, evidence, replay,
rollback, and degraded-mode survivability. \emph{DAD} argued that the
problem is institutional as much as technical: machine-executed
authority needs a joint control-plane authority, not just another
software program. \emph{From PDFs to Pull Requests} made the adjacent
point that governance cannot remain static prose. At machine tempo,
policy has to become executable, versioned, testable, and continuously
verified.

Recent NIST, DARPA, and Prometheus Flame signals do not replace that
argument.

They validate it.

\hypertarget{agents-are-not-apps}{%
\subsection{Agents are not apps}\label{agents-are-not-apps}}

This is the center of gravity.

Agents are not apps. They are delegated-authority systems.

An app usually executes inside a relatively fixed envelope. It takes
user input, processes it, and returns an output. An agent is different.
It plans. It selects actions. It uses tools. It traverses systems. It
can modify external state. It can continue operating after the initial
prompt. It can collaborate with other agents. It can create consequences
that outlive the interaction that started them.

Once that is true, the core problem changes.

The interesting question is no longer whether the model is impressive.

The question is whether the delegation is governable.

That requires a different architecture.

It requires identity for agents and sub-agents. It requires explicit
trust scopes rather than implied permission. It requires policy decision
points and policy enforcement points at the actual doing boundary:
runtime admission, retrieval, tool invocation, inter-agent messaging,
work-unit transitions, and release. It requires budgets, containment,
and revocation. It requires signed traces, provenance, replay, and
rollback. It requires environments that can be constrained and modified
rather than treated as an unbounded execution surface.

In other words, it requires a control plane.

That is the argument here. Not ``AI, but safer.'' Not ``better evals.''
Not ``more standards.''

A control plane for delegated authority.

\hypertarget{nist-is-starting-to-define-the-operational-surface}{%
\subsection{NIST is starting to define the operational
surface}\label{nist-is-starting-to-define-the-operational-surface}}

The most important thing in NIST's recent work is not any single
document. It is the shape of the portfolio.

The AI Agent Standards Initiative is explicitly about agents that can
act autonomously, function securely on behalf of users, and interoperate
through standards and open protocols. That is already a meaningful
shift. The target is no longer only model quality in isolation. The
target is trustworthy action in an ecosystem.

The January RFI on AI agent security makes the shift even clearer. It is
scoped to agent systems capable of affecting external state. It asks
about deployment environments, tool restrictions, continuous monitoring,
patching across the lifecycle, human oversight for consequential
actions, and the current state of ``undoes, rollbacks, or negations''
for unwanted trajectories. That is the right class of question. It means
the conversation is moving from ``is the model safe?'' to ``how do we
govern a system that can do things?''

The identity and authority concept work points in the same direction.
Identification, authorization, auditing, non-repudiation, and
prompt-injection controls are not accessory concerns. They are the
skeleton of accountable delegation.

The evaluation portfolio matters for the same reason. NIST AI 800-2
pushes toward validity, transparency, and reproducibility in automated
benchmark evaluations. NIST AI 800-3 pushes evaluators to formalize
assumptions and measurement targets statistically. NIST AI 800-4 shifts
attention from lab scoring to post-deployment monitoring. The MOU with
GSA ties evaluation more directly to procurement through USAi. The
OpenMined CRADA pushes on secure evaluation methods for settings where
data, models, or benchmarks cannot be openly shared.

Put differently, NIST is starting to build measurement science for trust
in operation.

That is good.

But measurement alone is not the answer.

Measurement needs a governed execution surface to measure.

\hypertarget{darpa-is-pushing-research-toward-the-same-problem}{%
\subsection{DARPA is pushing research toward the same
problem}\label{darpa-is-pushing-research-toward-the-same-problem}}

DARPA's AI Forward initiative lands in the same neighborhood from the
research side.

It frames the objective as trustworthy systems for national security
missions. It explicitly centers three thrusts: foundational theory, AI
engineering, and human-AI teaming. It says the point is to build systems
that work as intended in the real world, not just in the lab.

That is more important than it sounds.

For years, the AI discourse drifted toward a shallow equation: better
benchmark performance equals progress. DARPA is naming a deeper
requirement. Real systems have to operate reliably, interact
appropriately with people, and hold up in contested, messy environments.

Still, trustworthiness only becomes real when it has an engineering
surface.

A trustworthy system is not one that scores well on an eval card and
then disappears into runtime darkness. A trustworthy system is one whose
authority is bounded, whose actions are mediated, whose evidence is
retained, whose permissions can be reduced or revoked, whose behavior
can be supervised as a portfolio, and whose failure does not
automatically become a mission-scale event.

That is why trustworthy AI eventually becomes control-plane engineering.

\hypertarget{prometheus-flame-matters-because-acquisition-is-starting-to-buy-for-inheritance}{%
\subsection{Prometheus Flame matters because acquisition is starting to
buy for
inheritance}\label{prometheus-flame-matters-because-acquisition-is-starting-to-buy-for-inheritance}}

The Prometheus Flame draft matters less because of the deadline and more
because of what it revealed.

The Air Force did not frame the early phases as a competition for
another closed demo stack. The draft points instead toward open-source
AMS GRA and A-GRA implementations, permissive licensing, public
repositories and documentation, turn-key projects, CI/CD artifacts,
compliance reporting, and follow-on prototype and production OT
pathways. It also places real weight on usability, onboarding, standards
compliance, and risk rather than treating them as proposal garnish.

That is not normal demo procurement.

That is ecosystem formation.

The most important artifact in that kind of acquisition is not the thing
that flies once.

It is the thing other people can inherit.

A reference implementation. A starter kit. A conformance harness. A
public repo. A paved road. An example project that works out of the box.
A delivery system that others can extend without privileged access to
the original builder.

That logic sits naturally beside the argument in \emph{From PDFs to Pull
Requests}. If government wants an open, durable autonomy ecosystem,
standards cannot remain prose-only and compliance cannot remain
narrative-only. The substrate has to be executable: code, tests,
schemas, policy objects, evaluation gates, evidence events, reproducible
pipelines, and release discipline.

That is what it means to buy for inheritance.

\hypertarget{the-next-object-we-need-an-action-reversibility-profile}{%
\subsection{The next object we need: an Action Reversibility
Profile}\label{the-next-object-we-need-an-action-reversibility-profile}}

Previous writing on this site already argued for trust scopes, control
planes, policy gates, evidence, and rollback.

The extension I would add now is this: every serious agent deployment
should carry an \textbf{Action Reversibility Profile}.

The point is simple. Not all actions are equal, and our architectures
should stop pretending they are.

Some actions are safely reversible. Some are compensable but not
literally undoable. Some are effectively irreversible. Some should be
prohibited altogether.

That classification should not live in a policy memo on a shelf.

It should live in the control plane.

At minimum, an Action Reversibility Profile should classify actions into
four buckets:

\begin{enumerate}
\def\labelenumi{\arabic{enumi}.}
\tightlist
\item
  \textbf{Reversible}: actions that can be automatically undone with
  high confidence.
\item
  \textbf{Compensable}: actions that cannot be literally rolled back but
  can be corrected through a defined compensation path.
\item
  \textbf{Irreversible}: actions whose consequences cannot be
  practically reversed and therefore require tighter gating, stronger
  evidence, and usually explicit human authorization.
\item
  \textbf{Prohibited}: actions the system is never permitted to take.
\end{enumerate}

For each class, the profile should define:

\begin{itemize}
\tightlist
\item
  approval conditions
\item
  allowed tools and environments
\item
  evidence obligations
\item
  logging and provenance requirements
\item
  rollback or compensation mechanisms
\item
  timeout behavior
\item
  kill-switch behavior
\item
  update-handling rules
\item
  degraded-mode posture
\end{itemize}

This is where recent NIST signal becomes especially useful. If
government is already asking about deployment constraints, monitoring,
patching, human oversight, and the state of practice for undoes or
rollbacks, reversibility should become a first-class design property
rather than a hopeful runtime feature.

Reversibility is not just a safety pattern.

It is an operational tempo pattern.

The faster a system can act, the faster it must be able to stop,
contain, and recover.

\hypertarget{what-should-happen-next}{%
\subsection{What should happen next}\label{what-should-happen-next}}

If the current direction is real, the next moves should be concrete.

\hypertarget{standardize-authority-not-just-interfaces}{%
\subsubsection{1) Standardize authority, not just
interfaces}\label{standardize-authority-not-just-interfaces}}

Message schemas are not enough. The next standards layer has to carry
identity, trust scope, delegated permissions, policy hooks, and evidence
events.

\hypertarget{make-conformance-machine-checkable}{%
\subsubsection{2) Make conformance
machine-checkable}\label{make-conformance-machine-checkable}}

If a standard cannot be tested, it will eventually be narrated instead
of implemented. Conformance needs harnesses, reference scenarios,
negative tests, and machine-readable outputs.

\hypertarget{treat-rollback-as-a-first-class-design-property}{%
\subsubsection{3) Treat rollback as a first-class design
property}\label{treat-rollback-as-a-first-class-design-property}}

Rollback should shape action design, decomposition, and environment
mediation from the beginning, not appear as an afterthought after
deployment.

\hypertarget{collapse-procurement-evaluation-and-operations-into-one-evidence-loop}{%
\subsubsection{4) Collapse procurement, evaluation, and operations into
one evidence
loop}\label{collapse-procurement-evaluation-and-operations-into-one-evidence-loop}}

Procurement should not buy one thing, evaluation should not measure a
second thing, and operations should not trust a third thing. The same
evidence model should flow from source to test to deployment to
sustainment.

\hypertarget{fund-open-substrates-and-maintainers}{%
\subsubsection{5) Fund open substrates and
maintainers}\label{fund-open-substrates-and-maintainers}}

If the goal is a real ecosystem, government cannot only reward edge
demos. It has to fund the center: reference implementations, conformance
tooling, public starter kits, maintainers, and the delivery systems that
keep the whole thing governable over time.

\hypertarget{the-real-advantage}{%
\subsection{The real advantage}\label{the-real-advantage}}

The strategic advantage is not ``having AI.''

It is being able to field governed autonomy at mission tempo.

That means being able to delegate without surrendering authority. It
means being able to move faster without severing provenance. It means
being able to exploit open ecosystems without accepting open-ended risk.
It means many builders can contribute while identity, policy, evidence,
rollback, and containment remain coherent across the whole system.

That is the job of the control plane.

And that is why this moment matters.

NIST is starting to standardize toward the operational surface. DARPA is
starting to research toward trustworthy systems that work in the real
world. Defense acquisition is beginning to buy for open, inheritable
substrates rather than only polished integration theater.

Good.

Now the missing piece needs to be stated plainly.

The future will not be won by institutions with the most impressive
demos.

It will be won by institutions that learn how to govern delegated
authority.

That is the next layer.

That is the work.

\hypertarget{references}{%
\subsection{References}\label{references}}

\begin{enumerate}
\def\labelenumi{\arabic{enumi}.}
\tightlist
\item
  NIST CAISI, \emph{AI Agent Standards Initiative}.
  https://www.nist.gov/caisi/ai-agent-standards-initiative
\item
  NIST CAISI, \emph{CAISI Issues Request for Information About Securing
  AI Agent Systems}.
  https://www.nist.gov/news-events/news/2026/01/caisi-issues-request-information-about-securing-ai-agent-systems
\item
  NCCoE, \emph{New Concept Paper on Identity and Authority of Software
  Agents}.
  https://www.nccoe.nist.gov/news-insights/new-concept-paper-identity-and-authority-software-agents
\item
  NCCoE, \emph{Accelerating the Adoption of Software and AI Agent
  Identity and Authorization} (concept paper PDF).
  https://www.nccoe.nist.gov/sites/default/files/2026-02/accelerating-the-adoption-of-software-and-ai-agent-identity-and-authorization-concept-paper.pdf
\item
  NIST, \emph{Towards Best Practices for Automated Benchmark
  Evaluations} (NIST AI 800-2).
  https://www.nist.gov/news-events/news/2026/01/towards-best-practices-automated-benchmark-evaluations
\item
  NIST, \emph{NIST AI 800-2 Initial Public Draft} (PDF).
  https://nvlpubs.nist.gov/nistpubs/ai/NIST.AI.800-2.ipd.pdf
\item
  NIST, \emph{New Report: Expanding the AI Evaluation Toolbox with
  Statistical Models} (NIST AI 800-3).
  https://www.nist.gov/news-events/news/2026/02/new-report-expanding-ai-evaluation-toolbox-statistical-models
\item
  NIST, \emph{NIST AI 800-3} (PDF).
  https://nvlpubs.nist.gov/nistpubs/ai/NIST.AI.800-3.pdf
\item
  NIST, \emph{New Report: Challenges to the Monitoring of Deployed AI
  Systems} (NIST AI 800-4).
  https://www.nist.gov/news-events/news/2026/03/new-report-challenges-monitoring-deployed-ai-systems
\item
  NIST, \emph{NIST AI 800-4} (PDF).
  https://nvlpubs.nist.gov/nistpubs/ai/NIST.AI.800-4.pdf
\item
  NIST CAISI, \emph{CAISI signs MOU with GSA to boost AI evaluation
  science in federal procurement through USAi}.
  https://www.nist.gov/news-events/news/2026/03/caisi-signs-mou-gsa-boost-ai-evaluation-science-federal-procurement-through
\item
  NIST CAISI, \emph{CAISI signs CRADA with OpenMined to Enable Secure AI
  Evaluations}.
  https://www.nist.gov/news-events/news/2026/03/announcement-caisi-signs-crada-openmined-enable-secure-ai-evaluations
\item
  DARPA, \emph{AI Forward}.
  https://www.darpa.mil/research/programs/ai-forward
\item
  SAM.gov, \emph{Prometheus Flame - Request for Information} (Special
  Notice, FA8694-26-S-8888).
  https://sam.gov/opp/68a071f69a1c4ec883b8392c5ae040b1/view
\item
  SAM.gov API notice description (source for extracted draft text).
  https://api.sam.gov/prod/opportunities/v1/noticedesc?noticeid=68a071f69a1c4ec883b8392c5ae040b1
\end{enumerate}

\hypertarget{related-writing-on-this-site}{%
\subsection{Related writing on this
site}\label{related-writing-on-this-site}}

\begin{itemize}
\tightlist
\item
  \href{/writing/acp-ra/}{Agent Control Plane Reference Architecture
  (ACP-RA)}
\item
  \href{/writing/agentic-force-creation/}{From AI Force Multiplication
  to Force Creation}
\item
  \href{/writing/2026-02-12-stand-up-dad-joint-control-plane-authority/}{Stand
  Up Delegated Autonomy Directorate (DAD) as a Joint Control Plane
  Authority}
\item
  \href{/writing/code-as-policy/}{From PDFs to Pull Requests}
\end{itemize}

\end{document}
